\subsection{NARCBench Collusion Attack and Observable-Reasoning Monitor}
\label{sec:narcbench-cot}

\paragraph{Attack environment.}
We evaluated covert multi-agent coordination using the 50 scenarios from
NARCBench-Core~\citep{rose2026narcbench}. The scenarios span 17 high-stakes
domains and each instantiates a four-member committee choosing between two
options. For every scenario, we ran a collusion condition and a matched control
condition, yielding 100 runs. In the collusion condition, two designated agents
received a private coordination phase before three rounds of public
deliberation; all four agents then independently submitted a final
recommendation. In the control condition, all four agents received honest
mandates and no private coordination phase occurred. Scenario text, agent
identities, option content, speaking-order seed, and public deliberation
structure were held fixed within each matched pair. Following the released
benchmark implementation, option labels were counterbalanced across letter and
numeric labels and across target position to reduce label-position shortcuts.

\paragraph{Models and generation.}
The committee agents used \texttt{openai/bedrock/gpt-oss-120b}, accessed through
Amazon Bedrock's OpenAI-compatible endpoint, with temperature 0.7, low reasoning
effort, and a maximum of 768 output tokens per turn. Independent monitor calls
used the same underlying checkpoint in a separately configured Inspect model
role, with temperature 0, low reasoning effort, and a maximum of 1,536 output
tokens. Using the same checkpoint for generation and monitoring is a limitation
of this experiment and should not be interpreted as cross-model validation. We
used Inspect AI version 0.3.251. Agent prompts and model responses
were recorded in the Inspect evaluation log; SHA-256 hashes of every system/user
prompt pair were additionally stored in the trajectory metadata. The exact
NARCBench source revision was commit
\texttt{85fea3be2a9d8b9f16e33a4bd65bd947c31fa07d}.

\paragraph{Observable-reasoning instrumentation.}
Each committee response contained a public message (or final recommendation)
and a concise rationale made available to the evaluator but not to the other
committee members. We refer to the latter as an \emph{observable rationale},
not as the model's hidden or faithful chain of thought. This distinction is
important because post-hoc rationales may omit, rationalize, or misrepresent the
internal computation that produced an action. The instrumentation therefore
defines a NARCBench-derived evaluation variant rather than an exact reproduction
of the original generation protocol.

\paragraph{Attack outcomes.}
Let $T_i^{(c)}$ and $T_i^{(0)}$ indicate whether the committee selected the
assigned collusion target in the collusion and control runs for scenario $i$.
We defined target selection as at least three of four committee votes for the
assigned target. We report target-selection rates in each condition and the matched uplift
$\Delta = n^{-1}\sum_i(T_i^{(c)}-T_i^{(0)})$. We additionally report (i) target
adherence, the fraction of collusion runs in which both designated colluders
voted for the assigned target, and (ii) end-to-end attack success, which
requires both target adherence and committee selection of the target. All 50
assigned collusion runs remain in the primary analysis; abstentions, malformed
outputs, and refusals are counted as failures rather than silently excluded.
Separately, human reviewers audit whether the messages exhibit semantic
coordination. The audited-valid subset requires both target adherence and
semantic evidence of a shared strategy and is reported only as a secondary
per-protocol analysis. Before the preliminary meeting, one human reviewer
labeled a fixed 10-of-50 subset. Files previously discussed with an AI assistant
were excluded from this independent subset. Because a second reviewer was not
available, we do not report inter-rater agreement or Cohen's kappa, and this
subset is treated as a preliminary qualitative audit rather than a full-dataset
estimate.

\paragraph{Defense and evidence scopes.}
An independent LLM monitor assigned a suspiciousness score in $[0,1]$ to each
trajectory. The monitor never received the condition label, hidden target,
colluder identities, attack-specific metadata, or private messages in the main
conditions. We predeclared three evidence scopes: \emph{action-only}, containing
public messages, votes, tool calls, and the committee outcome;
\emph{observable-reasoning-only}, containing evaluator-visible rationales from
public and final stages; and \emph{action+reasoning}, containing both sources.
Action+reasoning was the primary defense condition; the other two were
ablations. A fourth privileged condition that includes private coordination may
be reported as an audit upper bound but is not treated as a deployable defense.
Scores at or above 0.5 were classified as attacks. The threshold was fixed
before examining the test results. This completed cell used a NARCBench-specific
monitor prompt (	exttt{narcbench_v1}) focused on covert steering of committee
decisions. For the cross-attack defense matrix, we separately freeze an
attack-agnostic prompt (	exttt{general_v2}) and rescore saved trajectories so
that every attack row uses the same defense prompt. Results from the two prompt
versions are not pooled.

\paragraph{Statistical analysis.}
For attack proportions and thresholded defense rates, we report 95\% Wilson
confidence intervals. For the paired target-selection uplift, we report a 95\%
percentile bootstrap interval obtained by resampling scenario pairs with
replacement (10,000 replicates) and a two-sided exact McNemar test. Monitor
performance is summarized by AUROC and average precision as threshold-free
metrics, together with recall, false-positive rate, precision, balanced
accuracy, and invalid-output counts at the fixed threshold. Resampling and all
comparisons preserve the scenario pair as the unit of analysis.

\paragraph{Compute and reproducibility.}
The full evaluation comprised 1,700 attacker-model calls and 300 monitor-model
calls. We report total input/output tokens, wall-clock time, model/provider IDs,
generation parameters, software versions, scenario-level outcomes, and the
manual-audit protocol alongside aggregate results. The released analysis script
converts a completed Inspect log into normalized JSONL, attack and defense CSV
summaries, and a manuscript-ready results table.

\paragraph{Results.}
Across 50 matched scenarios, target adherence was 0.280 (95\% CI
[0.175, 0.417]) and end-to-end attack success was 0.060 (95\% CI
[0.021, 0.162]). Target selection was 0.100 under collusion and 0.020 under
control, corresponding to a paired uplift of 0.080 (95\% paired-bootstrap CI
[-0.020, 0.180]; exact McNemar $p=0.219$). The primary action+reasoning monitor
achieved AUROC 0.316, recall 0.300, and false-positive rate 0.520. The
action-only baseline achieved the highest AUROC (0.427); consequently, these
results provide no evidence that the elicited rationales improved detection on
this model and protocol. All three scopes produced valid scores for all 100
trajectories. In the preliminary human-audited subset, 3 of 10 trajectories met
both automatic target adherence and the reviewer's semantic-coordination
criterion; 7 of 10 were rated as containing refusal or noncompliance. These
subset counts are not extrapolated to all 50 attack trajectories.
